\subsection{Fill and known-powers gadgets}
This section proves two mutually recursive statements.  The level-$l$ fill
gadget $A_{2^l}$ takes a compatible degree-$n$ pair, together with the monic
powers $H_2,\ldots,H_{2^l}$, and outputs a monic polynomial of degree
$n+2^{l+1}-1$.  Its decoder recovers the input pair and its
$\alpha$- and $\beta$-parameter blocks.  Feeding a simple
known-power pair into this gadget produces the decodable polynomial
$Q_{2^k-1}$ given $(H_2,\ldots,H_{2^{k-1}})$.  The proof of
\Cref{lem:fill-correctness} states the well-founded order explicitly, so this
mutual definition is not a circular decoding argument.  A depth-balanced
drop-in variant of the known-powers gadget, with the identical ledger and
logarithmic height, closes the section (\Cref{sec:peeled-Q}).
\Cref{fig:fill-slots} pictures the mutual recursion and its well-founded
decoding order.

\begin{figure}[H]
  \centering
  \resizebox{\textwidth}{!}{% One fill level, its pivot bands, and the well-founded mutual order.
\begin{tikzpicture}[x=1cm,y=1cm,>=Latex]
  \node[fp panel title] at (0,5.05) {(a) One fill level, $D=2^l$};
  \node[fp recursive, minimum width=27mm] (input) at (1.45,4.18)
    {input pair\\$(S_D^{(1)},S_D^{(2)})$};
  \node[fp known, minimum width=33mm] (heads) at (4.65,4.18)
    {heads $H_D+Q_{D/2-1}$\\and $H_D+\beta_D$};
  \node[fp second, minimum width=31mm] (mid) at (8.2,4.18)
    {multiply; add\\a fresh low $Q$-block};
  \node[fp active, minimum width=30mm] (lower) at (11.75,4.18)
    {intermediate pair\\at level $l-1$};
  \draw[fp flow] (input) -- (heads);
  \draw[fp flow] (heads) -- (mid);
  \draw[fp flow] (mid) -- (lower);
  \node[fp label, below=2mm of lower] {recurse, then output
    $P=(x+\beta_0)A^{(1)}+A^{(2)}$};

  \node[fp panel title] at (0,2.95) {(b) The three output pivot bands};
  \draw[fp decode] (0,2.5) -- node[fp label, above] {descending coefficient degree}
    (13.55,2.5);
  \node[fp active, minimum width=38mm, anchor=west] (betas) at (0,1.78)
    {$\beta_0,\ldots,\beta_D$\\unit pivots in the high band};
  \node[fp recursive, minimum width=47mm, anchor=west] (old) at (3.8,1.78)
    {input-pair pivots shifted by $2D-2$\\with their original slopes};
  \node[fp second, minimum width=46mm, anchor=west] (alphas) at (8.5,1.78)
    {$\alpha_{2D-3},\ldots,\alpha_0$\\unit pivots in the low band};
  \draw[fp dependence] (betas.south east) to[bend right=10] (old.south west);
  \draw[fp dependence] (old.south east) to[bend right=10] (alphas.south west);

  \node[draw=black!50, line width=.45pt, inner sep=2pt, font=\scriptsize,
    align=center, minimum width=105mm, minimum height=6mm] at (6.7,0.52)
    {well-founded proof order:
     $F_1\longrightarrow\mathsf Q_3\longrightarrow F_2
      \longrightarrow\mathsf Q_4\longrightarrow F_3\longrightarrow\cdots$,
     \quad$\mathsf Q_s$ means the decoder for $Q_{2^s-1}$};
  \node[fp label, anchor=west, text=fpBlue!80!black] at (0,-0.08)
    {A smaller known-powers gadget builds the heads; the completed fill level then proves the next larger gadget.};
\end{tikzpicture}
}
  \caption{Forward and inverse views of the fill construction.  One level
  multiplies the input pair by monic heads, adds a fresh low gadget, and
  delegates to the preceding fill level.  In the inverse direction its
  coefficients separate into three descending pivot bands: new top
  $\beta$-parameters, the shifted input-pair pivots, and new low
  $\alpha$-parameters.  The bottom line displays the well-founded order in
  which fill correctness and known-powers decodability are proved.
  The known-powers gadget assembled this way is the sequential baseline; the
  shipped schedules substitute the depth-balanced $Q^{\mathrm{peel}}$ of
  \Cref{sec:peeled-Q} at the identical ledger.}
  \label{fig:fill-slots}
\end{figure}

\begin{algorithm}[H]
    \caption{High-level fill construction}\label{alg:constr-fill}
    \begin{algorithmic}
        \If{$l = 1$}
            \Statex $A^{(1)}_{2}[\alpha_0, \alpha_1, \beta_2, \beta_1](S^{(1)}_2, (x, H_2)) = (H_2 + \beta_1) S^{(1)}_2 + \alpha_1$
            \Statex $A^{(2)}_{2}[\alpha_0, \alpha_1, \beta_2, \beta_1](S^{(2)}_2, (x, H_2)) = (H_2 + \beta_2) S^{(2)}_2 + \alpha_0$
        \ElsIf{$l = 2$}
            \Statex $S^{(1)}_{2} = (H_{4} + \beta_3) S^{(1)}_{4} + Q_{3}[\alpha_3, \alpha_4, \alpha_5](x, H_2)$
            \Statex $A^{(1)}_{4}[\alpha_0, \ldots, \alpha_5, \beta_3, \ldots, \beta_1](S^{(1)}_{4}, (x, H_2, H_4))$
            \Statex \hspace{\algorithmicindent}$= A^{(1)}_{2}[\alpha_0, \alpha_1, \beta_2, \beta_1](S^{(1)}_{2}, (x, H_2))$
            \Statex $S^{(2)}_{2} = (H_{4} + \beta_4) S^{(2)}_{4} + \alpha_2$
            \Statex $A^{(2)}_{4}[\alpha_0, \ldots, \alpha_5, \beta_4, \ldots, \beta_1](S^{(2)}_{4}, (x, H_2, H_{4}))$
            \Statex \hspace{\algorithmicindent}$= A^{(2)}_{2}[\alpha_0, \alpha_1, \beta_2, \beta_1](S^{(2)}_{2}, (x, H_2))$
        \Else
            \Statex $S^{(1)}_{2^{l - 1}} = (H_{2^{l}} + Q_{2^{l - 1} - 1}[\beta_{2^{l} - 1}, \ldots, \beta_{2^{l - 1} + 1}](x, H_2, \ldots, H_{2^{l - 2}})) S^{(1)}_{2^l}$
            \Statex \hspace{\algorithmicindent}$+\; Q_{2^{l} - 1}[\alpha_{2^{l} - 1}, \ldots, \alpha_{2^{l + 1} - 3}](x, H_2, \ldots, H_{2^{l - 1}})$
            \Statex $A^{(1)}_{2^l}[\alpha_0, \ldots, \alpha_{2^{l + 1} - 3}, \beta_{2^l - 1}, \ldots, \beta_{1}](S^{(1)}_{2^l}, (x, H_2, \ldots, H_{2^{l}}))$
            \Statex \hspace{\algorithmicindent}$= A^{(1)}_{2^{l - 1}}[\alpha_0, \ldots, \alpha_{2^l - 3}, \beta_{2^{l - 1} - 1}, \ldots, \beta_{1}](S^{(1)}_{2^{l - 1}}, (x, H_2, \ldots, H_{2^{l - 1}}))$
            \Statex $S^{(2)}_{2^{l - 1}} = (H_{2^{l}} + \beta_{2^{l}}) S^{(2)}_{2^l} + \alpha_{2^l - 2}$
            \Statex $A^{(2)}_{2^l}[\alpha_0, \ldots, \alpha_{2^{l + 1} - 3}, \beta_{2^l}, \ldots, \beta_{1}](S^{(2)}_{2^l}, (x, H_2, \ldots, H_{2^{l}}))$
            \Statex \hspace{\algorithmicindent}$= A^{(2)}_{2^{l - 1}}[\alpha_0, \ldots, \alpha_{2^l - 3}, \beta_{2^{l - 1}}, \ldots, \beta_{1}](S^{(2)}_{2^{l - 1}}, (x, H_2, \ldots, H_{2^{l - 1}}))$
        \EndIf
        \Statex $A_{2^l}[\alpha_0, \ldots, \alpha_{2^{l + 1} - 3}, \beta_{2^l}, \ldots, \beta_{0}](S^{(1)}_{2^l}, S^{(2)}_{2^l}, (x, H_2, \ldots, H_{2^{l}})) = (x + \beta_{0}) A^{(1)}_{2^l} + A^{(2)}_{2^l}$
    \end{algorithmic}
\end{algorithm}

The construction, \Cref{alg:constr-fill}, proceeds by recursively building intermediate pairs and finally outputting
$P=(x+\beta_0)A^{(1)}_{2^l}+A^{(2)}_{2^l}$.
We start by isolating this final step.

\begin{lemma}[$(x+\alpha)$-extraction]\label{lem:x-alpha-extraction}
    Let $(T^{(1)}, T^{(2)})$ be a compatible pair on
    $G \subseteq \idx{n}$ given $(B_i)_{i \in [t]}$, with
    $\deg T^{(1)}=\deg T^{(2)}=n$.
    Define $P[\alpha](x)=(x+\alpha)\,T^{(1)}(x)+T^{(2)}(x)$.
    If $G\subseteq\rng n$, then: (1) $\alpha$ is extractable from $P$
    given $(B_i)_{i\in[t]}$; and (2) $T^{(1)}$ and $T^{(2)}$ are
    derivable from $\bigl(P,(B_i)_{i\in[t]}\bigr)$.
    More generally, for an arbitrary $G\subseteq\idx n$, conclusion~(2)
    still holds whenever $\alpha$ has already been derived.
\end{lemma}
\begin{proof}
    First suppose $G \subseteq \rng{n}$.  Compatibility for $j=n-1$
    then expresses $\coeff{T^{(1)}}{n-1}$ as a polynomial in the
    $(B_i)$-coefficients (the tuple
    $\bigl(\coeff{xT^{(1)}+T^{(2)}}{i}\bigr)_{i\in G,\, i\ge n}$ is
    empty).
    Moreover, since $T^{(1)}$ and $T^{(2)}$ are monic of degree $n$,
    \[
        \coeff{P}{n}
            = \coeff{(x+\alpha)T^{(1)}}{n} + \coeff{T^{(2)}}{n}
            = \alpha\,\coeff{T^{(1)}}{n} + \coeff{T^{(1)}}{n-1} + 1
            = \alpha + \coeff{T^{(1)}}{n-1} + 1.
    \]
    Hence $\alpha=\coeff{P}{n}-\coeff{T^{(1)}}{n-1}-1$ is extractable from $P$ given $(B_i)_{i\in[t]}$.

    For the second conclusion, no restriction beyond $G\subseteq\idx n$
    is needed once $\alpha$ is known.  Define
    $\Phi(x)\coloneqq xT^{(1)}(x)+T^{(2)}(x)$. Since
    $P=\Phi+\alpha T^{(1)}$, we have
    \[
        \coeff{\Phi}{k}=\coeff{P}{k}-\alpha\,\coeff{T^{(1)}}{k}\qquad(k\ge 0).
    \]
    We recover the coefficients of $T^{(1)}$ and $T^{(2)}$ by descending induction on $j=n-1,n-2,\ldots,0$.
    The leading coefficients in degree $n$ are already known to be one.
    The induction hypothesis is that $\coeff{T^{(1)}}{r}$ and
    $\coeff{T^{(2)}}{r}$ are known for all $r>j$, and thus
    $\coeff{\Phi}{i}=\coeff{P}{i}-\alpha\coeff{T^{(1)}}{i}$ is known for
    all $i\in G$ with $i>j$.  This includes the possible top row $i=n$.
    By compatibility, $\coeff{T^{(1)}}{j}$ is a polynomial in $\bigl(\coeff{\Phi}{i}\bigr)_{i\in G,\,i\ge j+1}$ and the $(B_i)$-coefficients; since $i\ge j+1$ implies $i>j$, these $\Phi$-coefficients are already known by the induction hypothesis.
    Once $\coeff{T^{(1)}}{j}$ is known, we can compute $\coeff{\Phi}{j}=\coeff{P}{j}-\alpha\,\coeff{T^{(1)}}{j}$, and then compatibility expresses $\coeff{T^{(2)}}{j}$ as a polynomial in $\bigl(\coeff{\Phi}{i}\bigr)_{i\in G,\,i\ge j}$ and the $(B_i)$-coefficients.
    This completes the induction and shows that $T^{(1)}$ and $T^{(2)}$ are derivable from $\bigl(P,(B_i)_{i\in[t]}\bigr)$.
\end{proof}

\begin{lemma}[Aux-head compatibility for degree $2$]\label{lem:compatible-auxhead-deg2}
    Let $H_2\in\mathbb{F}[x]$ be monic of degree $2$, and let $b_1,b_2$ be scalars.
    Define $F_1\coloneqq H_2+b_1$ and $F_2\coloneqq H_2+b_2$.
    Then $(F_1,F_2)$ is a compatible pair on $\{0,1\}$ given $(H_2)$.
\end{lemma}
\begin{proof}
    Define $\Phi\coloneqq xF_1+F_2$.
    Since $F_1$ and $F_2$ are monic of degree $2$, their leading coefficients are known.
    Also $\coeff{F_1}{1}=\coeff{F_2}{1}=\coeff{H_2}{1}$ is derivable from $H_2$.

    Finally, using that $\coeff{\Phi}{0}=\coeff{F_2}{0}$ and $\coeff{\Phi}{1}=\coeff{F_1}{0}+\coeff{F_2}{1}$, we can derive
    \[
        \coeff{F_2}{0}=\coeff{\Phi}{0},
        \qquad
        \coeff{F_1}{0}=\coeff{\Phi}{1}-\coeff{H_2}{1}.
    \]
    The unknown constant coefficient of $F_1$ is therefore read from row~$1$,
    while that of $F_2$ is read from row~$0$; their degree-one coefficients
    and all higher coefficients are supplied by $H_2$.  Hence the first
    component uses only rows at least $j+1$ and the second only rows at least
    $j$, which witnesses compatibility on $\{0,1\}$ given $(H_2)$.
\end{proof}

\begin{lemma}[A trivial compatible pair for adding constants]\label{lem:compatible-monic-plus-constants}
    Let $n\ge 2$ and let $u,v$ be scalars.
    Then the pair $\bigl(x^n+u,\;x^n+v\bigr)$ is compatible on $\{0,1\}$.
\end{lemma}
\begin{proof}
    Define $\Phi\coloneqq x(x^n+u)+(x^n+v)=x^{n+1}+x^n+ux+v$.
    Since $n\ge 2$, we have $\coeff{\Phi}{0}=v$ and $\coeff{\Phi}{1}=u$.
    Thus the only unknown coefficient of the first component, $u$ in degree
    zero, is read from row $1$, and the only unknown coefficient of the second
    component, $v$ in degree zero, is read from row $0$; every other
    coefficient is fixed.  These are exactly the two compatibility cutoffs.
\end{proof}

\noindent
In later constructions we repeatedly meet pairs of the form $(H+Q,\;H+b)$ where $H$ is a known monic
polynomial, $Q$ is an \emph{unknown} lower-degree monic polynomial, and $b$ is an unknown scalar shift.
The next lemma records that such a pair is compatible on the low-degree window that contains all
coefficients of $Q$: the map $\Phi=x(H+Q)+(H+b)$ is block-triangular, so one can peel off $Q$ (and $b$)
from low-degree coefficients of $\Phi$ given $H$.

\begin{lemma}[A block-triangular compatible pair]\label{lem:compatible-aux-add-left}
    Let $H\in\mathbb{F}[x]$ be monic of degree $m\ge 1$, let $Q\in\mathbb{F}[x]$ be monic of degree $q<m$, and let $b$ be a scalar.
    Define $P^{(1)}\coloneqq H+Q$ and $P^{(2)}\coloneqq H+b$.
    Then $(P^{(1)},P^{(2)})$ is a compatible pair on $\idx{q}$ given $(H)$.
\end{lemma}
\begin{proof}
    Define $\Phi\coloneqq xP^{(1)}+P^{(2)}$ and note that
    \[
      \Phi = x(H+Q) + (H+b) = (xH+H) + (xQ+b).
    \]
    Since $H$ is given, the polynomial $xH+H$ is derivable.
    Therefore the polynomial
    \[
      \Xi \coloneqq \Phi-(xH+H)=xQ+b
    \]
    is derivable from $\Phi$ given $(H)$.
    In particular, $b=\coeff{\Xi}{0}$ is derivable.

    Next, for $1\le i\le q$ we have $\coeff{\Xi}{i}=\coeff{Q}{i-1}$, so the coefficients $\coeff{Q}{0},\ldots,\coeff{Q}{q-1}$ are derivable from $\bigl(\coeff{\Phi}{i}\bigr)_{i\in\idx{q}}$ given $(H)$.
    Together with monicity $\coeff{Q}{q}=1$, this derives $Q$, hence also $P^{(1)}=H+Q$ and $P^{(2)}=H+b$.

    It remains to verify the coefficient-dependency condition in the definition of compatibility.
    Since $q<m$, we have $\coeff{P^{(1)}}{j}=\coeff{P^{(2)}}{j}=\coeff{H}{j}$ for all $j>q$, and these coefficients are derivable from $(H)$ alone.
    For $j\le q$, each coefficient of $Q$ (and hence of $P^{(1)}$) is derivable from the coefficients $\bigl(\coeff{\Phi}{i}\bigr)_{i\in\idx{q},\, i\ge j+1}$ and $(H)$ by the explicit reconstruction above (since $\coeff{Q}{j}$ depends only on $\coeff{\Xi}{j+1}$ when $j<q$, and $\coeff{Q}{q}=1$).
    Similarly, $\coeff{P^{(2)}}{0}$ is derivable from $\coeff{\Phi}{0}$ and $(H)$ via $b$, and $\coeff{P^{(2)}}{j}=\coeff{H}{j}$ for $j\ge 1$.
    This gives the required coefficient polynomials in the definition of compatibility on $\idx{q}$ given $(H)$.
\end{proof}

\begin{lemma}[Padding a low pair]\label{lem:compatible-low-padding}
Let $N\ge r\ge1$, let $Q$ be monic of degree $r-1$, and let $b$ be a
scalar.  Then
\[
  (x^N+Q,\ x^N+b)
\]
is compatible on $\rng r$.
\end{lemma}
\begin{proof}
After subtracting the known terms $x^{N+1}+x^N$ from its combined
polynomial, the window $\rng r$ is exactly the corresponding window of
$xQ+b$.  Degree zero gives $b$, and degree $j+1$ gives $[x^j]Q$ for
$0\le j<r-1$; the remaining coefficient of $Q$ is its known leading~$1$.
Moreover $[x^j]Q$ is read in degree $j+1$ and $b$ in degree zero, so these
formulas satisfy the two compatibility cutoffs.
\end{proof}

The known-powers construction below calls the fill gadget, while higher fill
levels call smaller known-powers gadgets.  The bases $Q_1,Q_3$ and the
well-founded order in \Cref{lem:fill-correctness} resolve this mutual
recursion.

\begin{algorithm}[H]
    \caption{Known powers construction}\label{alg:constr-known-2n-1}
    \begin{algorithmic}
        \If{$k = 1$}
            \Statex $Q_{1}[\alpha_0](x)=x+\alpha_0$
        \ElsIf{$k = 2$}
            \Statex $Q_{3}[\alpha_0, \alpha_1, \alpha_2](x, H_2) = (x + \alpha_2)(H_2 + \alpha_1) + \alpha_0$
        \ElsIf{$k = 3$}
            \Statex $S^{(1)}_{2} = H_{4} + \alpha_{3}$
            \Statex $S^{(2)}_{2} = H_{4} + \alpha_{2}$
            \Statex $Q_{7}[\alpha_0, \ldots, \alpha_{6}](x, H_2, H_{4}) = A_{2}[\alpha_0, \alpha_1, \alpha_4, \alpha_5, \alpha_6](S^{(1)}_{2}, S^{(2)}_{2}, (x, H_2))$
        \Else
            \Statex $S^{(1)}_{2^{k - 2}} = H_{2^{k - 1}} + Q_{2^{k - 2} - 1}[\alpha_{2^{k - 1} - 1}, \ldots, \alpha_{2^{k - 1} + 2^{k - 2} - 3}](x, H_2, \ldots, H_{2^{k - 3}})$
            \Statex $S^{(2)}_{2^{k - 2}} = H_{2^{k - 1}} + \alpha_{2^{k - 1} - 2}$
            \Statex $Q_{2^k - 1}[\alpha_0, \ldots, \alpha_{2^k - 2}](x, H_2, \ldots, H_{2^{k - 1}})$
            \Statex \hspace{\algorithmicindent}$= A_{2^{k - 2}}[\alpha_0, \ldots, \alpha_{2^{k - 1} - 3}, \alpha_{2^{k - 1} + 2^{k - 2} - 2}, \ldots, \alpha_{2^k - 2}](S^{(1)}_{2^{k - 2}}, S^{(2)}_{2^{k - 2}}, (x, H_2, \ldots, H_{2^{k - 2}}))$
        \EndIf
    \end{algorithmic}
\end{algorithm}

The separate $k=3$ line is an addition-count optimization, not a distinct
decoding principle.  The uniform recursive formula could instead take
$S^{(1)}_2=H_4+Q_1=H_4+x+\alpha_3$; the same fill decoder would apply and the
multiplication count would remain three, but it would use one additional
polynomial addition.  We retain the displayed peephole because the exact
addition count below records the optimized value~eight.  The schedules
certified in Lean use the depth-balanced drop-in $Q^{\mathrm{peel}}$ of
\Cref{sec:peeled-Q} (\cref{eq:peeled-Q}) in place of every instance of this
construction; it has the same parameter interface and gate ledger, and its
own decoder (\Cref{lem:peeled-Q-decodable}).

We first record the two direct decoder bases.

\begin{lemma}
    $Q_{1}[\alpha_0](x)$ is decodable.
\end{lemma}
\begin{proof}
    By definition $Q_1(x)=x+\alpha_0$, so $\alpha_0=[x^0]Q_1$ is extractable from $Q_1$.
\end{proof}

\begin{lemma}
    $Q_{3}[\alpha_0, \alpha_1, \alpha_2](x, H_2)$ is decodable given $H_2$.
\end{lemma}
\begin{proof}
    By definition, $Q_3(x,H_2)=(x+\alpha_2)(H_2+\alpha_1)+\alpha_0$ with $H_2$ monic of degree 2.
    Writing $H_2(x)=x^2 + h_1 x + h_0$, a direct expansion gives
    \begin{align}
        [x^2]Q_3 &= h_1 + \alpha_2, \\
        [x^1]Q_3 &= h_0 + \alpha_1 + h_1\alpha_2, \\
        [x^0]Q_3 &= \alpha_0 + h_0\alpha_2 + \alpha_1\alpha_2.
    \end{align}
    Hence
    \begin{align}
        \alpha_2 &= [x^2]Q_3 - [x^1]H_2, \\
        \alpha_1 &= [x^1]Q_3 - [x^0]H_2 - \alpha_2\,[x^1]H_2, \\
        \alpha_0 &= [x^0]Q_3 - \alpha_2\bigl([x^0]H_2 + \alpha_1\bigr),
    \end{align}
    which are polynomials in the coefficients of $Q_3$ and $H_2$. Thus $Q_3$ is decodable given $H_2$.
\end{proof}

We now prove both decoder statements in their actual dependency order.

\begin{lemma}\label{lem:fill-correctness}
    Let $L\ge 1$ and let $H_2, \ldots, H_{2^L}$ be a sequence of monic polynomials with $\deg H_{2^i} = 2^i$ for $1 \le i \le L$.
    Fix $l$ with $1\le l\le L$.
    (We allow extra known powers $H_{2^{l+1}},\ldots,H_{2^L}$ since later arguments may carry more auxiliary data than the particular fill level $l$ needs.)
    Assume that $(S^{(1)}_{2^l}, S^{(2)}_{2^l})$ are a compatible pair on $G \subseteq \rng{n - 2^l}$ given $(H_2, \ldots, H_{2^L})$ with $n = \deg S^{(1)}_{2^l} = \deg S^{(2)}_{2^l} \ge 2^l$.
    Apply the level-$l$ fill with fresh parameter blocks as follows:
    \[
      \begin{aligned}
        P\coloneqq{}&A_{2^l}
          [\alpha_0,\ldots,\alpha_{2^{l+1}-3},
           \beta_{2^l},\ldots,\beta_0]\\
        &\bigl(S^{(1)}_{2^l},S^{(2)}_{2^l},
               (x,H_2,\ldots,H_{2^l})\bigr).
      \end{aligned}
    \]
    Then the monic polynomial $P$
    satisfies that $\alpha_0, \ldots, \alpha_{2^{l + 1} - 3}, \beta_{2^l}, \ldots, \beta_{0}$ are extractable from $P$ given $(H_2, \ldots, H_{2^L})$, and that $S^{(1)}_{2^l}$ and $S^{(2)}_{2^l}$ are derivable from $(P, H_2, \ldots, H_{2^L})$.

    In particular, $Q_{2^k -1}[\alpha_0, \ldots, \alpha_{2^k - 2}]$ is decodable given $(H_2,\ldots,H_{2^{k-1}})$ for $k \ge 2$.
\end{lemma}
\begin{proof}
    The fill and known-powers assertions are mutually recursive, so we make
    the well-founded order explicit.  Let $F_l$ denote the fill assertion at
    level $l$ (uniformly for every $L\ge l$ and every permitted input pair),
    and let $\mathsf Q_s$ denote decodability of $Q_{2^s-1}$ from its known
    powers.  The displayed formulas above prove $\mathsf Q_1$ and
    $\mathsf Q_2$ directly.  We prove,
    for increasing $l$,
    \[
       (F_1,\ldots,F_{l-1};\ \mathsf Q_1,\ldots,\mathsf Q_l)\Longrightarrow F_l,
       \qquad
       (F_l;\ \mathsf Q_1,\ldots,\mathsf Q_l)\Longrightarrow \mathsf Q_{l+2}.
       \tag{fill-order}
    \]
    For $l=1$ the second implication is the special $Q_7$ call; for
    $l\ge2$ it is the ordinary outer fill call.  Thus the logical order is
    $F_1,\mathsf Q_3,F_2,\mathsf Q_4,F_3,\mathsf Q_5,\ldots$, and no
    construction is used before its decoder has been proved.  We now prove
    the first implication.

    Base case $l=1$: The output is
    \[
        P=(x+\beta_0)\bigl((H_2+\beta_1)S^{(1)}_2+\alpha_1\bigr)+\bigl((H_2+\beta_2)S^{(2)}_2+\alpha_0\bigr).
    \]
    We claim that $(A^{(1)}_2,A^{(2)}_2)$ is a compatible pair on
    \[
        G' \coloneqq (2+G)\cup(n+\{0,1\})\cup\{0,1\}\subseteq\rng{n+2}
    \]
    given $(H_2,\ldots,H_{2^L})$.
    Indeed, by \Cref{lem:compatible-auxhead-deg2} the degree-$2$ pair $\bigl(H_2+\beta_1,\;H_2+\beta_2\bigr)$ is compatible on $\{0,1\}$ given $(H_2)$.
    Since $G\subseteq\rng{n-2}$, the shifted windows $(2+G)$ and $(n+\{0,1\})$ are disjoint, so the Multiplicativity lemma implies that $\bigl((H_2+\beta_1)S^{(1)}_2,\;(H_2+\beta_2)S^{(2)}_2\bigr)$ is compatible on $(2+G)\cup(n+\{0,1\})$ given $(H_2,\ldots,H_{2^L})$.
    Moreover, by \Cref{lem:compatible-monic-plus-constants} (applied with $n\gets n+2$), the pair $(x^{n+2}+\alpha_1,\;x^{n+2}+\alpha_0)$ is compatible on $\{0,1\}$, and this window is disjoint from $(2+G)\cup(n+\{0,1\})$.
    Therefore, applying the Additivity lemma in the equal-degree case yields that $\bigl(A^{(1)}_2,A^{(2)}_2\bigr)=\bigl((H_2+\beta_1)S^{(1)}_2+\alpha_1,\;(H_2+\beta_2)S^{(2)}_2+\alpha_0\bigr)$ is compatible on $G'$ given $(H_2,\ldots,H_{2^L})$.
    Now we may apply \Cref{lem:x-alpha-extraction} to $P=(x+\beta_0)A^{(1)}_2+A^{(2)}_2$ to extract $\beta_0$ and derive $A^{(1)}_2,A^{(2)}_2$.

    Since $(S^{(1)}_2,S^{(2)}_2)$ is compatible on $G\subseteq\rng{n-2}$ given $(H_2,\ldots,H_{2^L})$, its coefficients of degrees $n-1$ and $n-2$ are derivable from the given data alone.
    Comparing the coefficient of $x^{n}$ in $A^{(1)}_2=(H_2+\beta_1)S^{(1)}_2+\alpha_1$ then extracts $\beta_1$, and similarly extracts $\beta_2$ from $A^{(2)}_2=(H_2+\beta_2)S^{(2)}_2+\alpha_0$.
    Since $H_2+\beta_1$ and $H_2+\beta_2$ are monic, \Cref{lem:monic-division} applied to
    \(
      A^{(1)}_2=(H_2+\beta_1)S^{(1)}_2+\alpha_1
    \)
    and
    \(
      A^{(2)}_2=(H_2+\beta_2)S^{(2)}_2+\alpha_0
    \)
    derives $S^{(1)}_2,S^{(2)}_2$ and the remainders $\alpha_1,\alpha_0$, completing the base case.

    Induction step: Assume the claim holds for all smaller levels, and fix $l\ge 2$.
    The recursive equalities in \Cref{alg:constr-fill} show that $P$ is obtained by applying the level-$(l-1)$ fill construction to the intermediate pair $(S^{(1)}_{2^{l-1}},S^{(2)}_{2^{l-1}})$ (with parameters $\alpha_0,\ldots,\alpha_{2^l-3}$ and $\beta_{2^{l-1}},\ldots,\beta_0$).

    We first check that this intermediate pair satisfies the assumptions needed to apply the induction hypothesis at level $l-1$.
    Write $n\coloneqq\deg S^{(1)}_{2^l}=\deg S^{(2)}_{2^l}$, so $\deg S^{(1)}_{2^{l-1}}=\deg S^{(2)}_{2^{l-1}}=n+2^l$.

    If $l=2$, \Cref{alg:constr-fill} defines
    \[
      S^{(1)}_{2}=(H_{4}+\beta_3)\,S^{(1)}_{4}+Q_{3}[\alpha_3,\alpha_4,\alpha_5](x,H_2),
      \qquad
      S^{(2)}_{2}=(H_{4}+\beta_4)\,S^{(2)}_{4}+\alpha_2.
    \]
    Since $H_4$ is given, subtracting $xH_4+H_4$ from $x(H_4+\beta_3)+(H_4+\beta_4)$ isolates $\beta_3x+\beta_4$, so $\bigl(H_4+\beta_3,\;H_4+\beta_4\bigr)$ is compatible on $\{0,1\}$ given $(H_4)$.
    Together with compatibility of $(S^{(1)}_{4},S^{(2)}_{4})$ on $G\subseteq\rng{n-4}$, the Multiplicativity lemma shows that the product pair
    \(
      \bigl((H_{4}+\beta_3)S^{(1)}_{4},\;(H_{4}+\beta_4)S^{(2)}_{4}\bigr)
    \)
    is compatible on $(4+G)\cup(n+\{0,1\})$.
    Put $N=n+4$.  By \Cref{lem:compatible-low-padding}, the padded low pair
    $(x^N+Q_3,x^N+\alpha_2)$ is compatible on $\rng4$.  Its window is
    disjoint from the product window, so Additivity (with the compensating
    $-x^N$ in its equal-degree case) shows that
    $(S^{(1)}_2,S^{(2)}_2)$ is compatible on
    \[
      \rng4\cup(4+G)\cup(n+\{0,1\})\subseteq\rng{(n+4)-2}.
    \]

    Now assume $l\ge 3$.
    Let
    \[
      \begin{aligned}
      Q_{\mathrm{low}}\coloneqq{}&
        Q_{2^{l-1}-1}[\beta_{2^l-1},\ldots,\beta_{2^{l-1}+1}]
          (x,H_2,\ldots,H_{2^{l-2}}),\\
      Q_{\mathrm{high}}\coloneqq{}&
        Q_{2^l-1}[\alpha_{2^l-1},\ldots,\alpha_{2^{l+1}-3}]
          (x,H_2,\ldots,H_{2^{l-1}}).
      \end{aligned}
    \]
    By \Cref{alg:constr-fill}, the intermediate pair is defined by
    \[
      S^{(1)}_{2^{l-1}}=(H_{2^l}+Q_{\mathrm{low}})\,S^{(1)}_{2^l}+Q_{\mathrm{high}},
      \qquad
      S^{(2)}_{2^{l-1}}=(H_{2^l}+\beta_{2^l})\,S^{(2)}_{2^l}+\alpha_{2^l-2}.
    \]
    By \Cref{lem:compatible-aux-add-left}, the pair $\bigl(H_{2^l}+Q_{\mathrm{low}},\;H_{2^l}+\beta_{2^l}\bigr)$ is compatible on $\idx{2^{l-1}-1}$ given $(H_{2^l})$.
    Since $(S^{(1)}_{2^l},S^{(2)}_{2^l})$ is compatible on $G$ and $G\subseteq\rng{n-2^l}$, we have
    \[
      (n+\idx{2^{l-1}-1})\cap(2^l+G)=\emptyset,
    \]
    and therefore the Multiplicativity lemma implies that the product pair
    \[
      \left(
      \begin{aligned}
        &(H_{2^l}+Q_{\mathrm{low}})S^{(1)}_{2^l},\\[-2pt]
        &(H_{2^l}+\beta_{2^l})S^{(2)}_{2^l}
      \end{aligned}
      \right)
    \]
    is compatible on $(2^l+G)\cup(n+\idx{2^{l-1}-1})$ given $(H_2,\ldots,H_{2^L})$.
    Put $N=n+2^l$.  By \Cref{lem:compatible-low-padding}, the padded pair
    $(x^N+Q_{\mathrm{high}},x^N+\alpha_{2^l-2})$ is compatible on
    $\rng{2^l}$.  This is disjoint from the displayed product window, so a
    further application of Additivity shows that the intermediate pair is
    compatible on
    \[
      G_{\mathrm{int}}\coloneqq \rng{2^l}\cup(2^l+G)\cup(n+\idx{2^{l-1}-1})
      \subseteq \rng{(n+2^l)-2^{l-1}}
    \]
    given $(H_2,\ldots,H_{2^L})$.

    By the induction hypothesis at level $l-1$, the parameters
    $\alpha_0,\ldots,\alpha_{2^l-3}$ and
    $\beta_{2^{l-1}},\ldots,\beta_0$ are extractable from $P$ given the
    powers.  The same hypothesis derives
    $(S^{(1)}_{2^{l-1}},S^{(2)}_{2^{l-1}})$ from $P$ and those powers.

    If $l=2$, we now recover $(S^{(1)}_{4},S^{(2)}_{4})$ and the remaining parameters as follows.
    Since $(S^{(1)}_{4},S^{(2)}_{4})$ is compatible on $G\subseteq\rng{n-4}$, its top four coefficients are derivable from the given power data alone.
    Comparing the coefficient of $x^{n}$ in $S^{(1)}_{2}=(H_{4}+\beta_3)S^{(1)}_{4}+Q_3$ then extracts $\beta_3$ (the $Q_3$ term does not contribute to degree $n$), and similarly extracts $\beta_4$ from $S^{(2)}_{2}=(H_{4}+\beta_4)S^{(2)}_{4}+\alpha_2$.
    Applying \Cref{lem:monic-division} to $S^{(1)}_{2}$ and $H_4+\beta_3$ derives $S^{(1)}_{4}$ and the remainder $Q_3$, and applying it to $S^{(2)}_{2}$ and $H_4+\beta_4$ derives $S^{(2)}_{4}$ and the remainder $\alpha_2$.
    Finally, the earlier $Q_3$ lemma decodes $Q_3$ to extract $\alpha_3,\alpha_4,\alpha_5$.
    This completes the induction step for $l=2$.

    Now assume $l\ge 3$.
    It remains to recover the two blocks
    \[
      \beta_{2^l},\beta_{2^l-1},\ldots,\beta_{2^{l-1}+1}
      \quad\text{and}\quad
      \alpha_{2^l-2},\alpha_{2^l-1},\ldots,\alpha_{2^{l+1}-3},
    \]
    and to derive $(S^{(1)}_{2^l},S^{(2)}_{2^l})$.
    Define
    \[
      \Phi\coloneqq xS^{(1)}_{2^{l-1}}+S^{(2)}_{2^{l-1}}.
    \]
    Using the displayed definitions and $\deg Q_{\mathrm{high}}<2^l$, the polynomial $xQ_{\mathrm{high}}+\alpha_{2^l-2}$ contributes only to degrees $\le 2^l$, and its contribution in degree $2^l$ is the known constant $1$ (since $Q_{\mathrm{high}}$ is monic of degree $2^l-1$).
    Therefore, from $\Phi$ we can derive the combined product polynomial
    \[
      \begin{aligned}
      \Phi_{\mathrm{prod}}\coloneqq{}&\Phi-x^{2^l}\\
        ={}&x\bigl((H_{2^l}+Q_{\mathrm{low}})S^{(1)}_{2^l}\bigr)
           +(H_{2^l}+\beta_{2^l})S^{(2)}_{2^l}\\
        &\quad +(\text{terms supported in degrees }<2^l).
      \end{aligned}
    \]
    In particular, the coefficients of $\Phi_{\mathrm{prod}}$ in the window $(2^l+G)\cup(n+\idx{2^{l-1}-1})$ agree with those of
    \(x\bigl((H_{2^l}+Q_{\mathrm{low}})S^{(1)}_{2^l}\bigr) + (H_{2^l}+\beta_{2^l})S^{(2)}_{2^l}\).
    Since the corresponding product pair is compatible on this window, we can derive
    \(
      (H_{2^l}+Q_{\mathrm{low}})S^{(1)}_{2^l}
    \)
    and
    \(
      (H_{2^l}+\beta_{2^l})S^{(2)}_{2^l}
    \),
    and then the Multiplicativity lemma derives the factors $(H_{2^l}+Q_{\mathrm{low}},H_{2^l}+\beta_{2^l})$ and $(S^{(1)}_{2^l},S^{(2)}_{2^l})$ from these products given $(H_2,\ldots,H_{2^L})$.

    With these polynomials derived, we can extract $\beta_{2^l}$ from the constant term of $H_{2^l}+\beta_{2^l}$.
    Also, applying \Cref{lem:monic-division} to $S^{(1)}_{2^{l-1}}$ and the monic divisor $H_{2^l}+Q_{\mathrm{low}}$ derives the remainder $Q_{\mathrm{high}}$, and applying it to $S^{(2)}_{2^{l-1}}$ and $H_{2^l}+\beta_{2^l}$ derives the remainder $\alpha_{2^l-2}$.

    Finally, to extract the remaining $\beta$- and $\alpha$-blocks, observe that $Q_{\mathrm{low}}$ and $Q_{\mathrm{high}}$ are the known-powers gadgets with exponents $s=l-1$ and $s=l$, respectively.  Both belong to the already-proved list $\mathsf Q_1,\ldots,\mathsf Q_l$ in (fill-order).  Hence they are decodable given the relevant known powers.
    Hence the remaining $\beta$-block is extractable from $Q_{\mathrm{low}}$, and the remaining $\alpha$-block is extractable from $Q_{\mathrm{high}}$.
    This proves the first implication, namely $F_l$.

    It remains to verify the second implication in (fill-order).  For $l=1$,
    the $Q_7$ call takes as input
    $(H_4+\alpha_3,H_4+\alpha_2)$, which is compatible on $\{0,1\}$ given
    $H_4$, because subtracting $xH_4+H_4$ from its combined polynomial
    leaves $\alpha_3x+\alpha_2$.  Thus $F_1$ proves $\mathsf Q_3$.  For $l\ge2$,
    put $k=l+2$ in \Cref{alg:constr-known-2n-1}.  Its level-$l$ outer fill
    call has input pair
    \[
      \bigl(H_{2^{l+1}}+Q_{2^l-1},
            H_{2^{l+1}}+\alpha_{2^{l+1}-2}\bigr).
    \]
    This pair is compatible on
    $\idx{2^l-1}=\rng{2^l}$ given $H_{2^{l+1}}$ by
    \Cref{lem:compatible-aux-add-left}; the inner $Q_{2^l-1}$ is the
    already-proved assertion $\mathsf Q_l$.  Applying $F_l$, then subtracting
    $H_{2^{l+1}}$ from the recovered input pair and invoking the decoder for
    $Q_{2^l-1}$, proves $\mathsf Q_{l+2}$.  This closes the simultaneous induction and
    proves the ``in particular'' assertion for every $k\ge2$.
\end{proof}

\begin{lemma}[Coefficient pivots of the known-powers gadgets]
\label{lem:Q-unitriangular}
Let $q=2^s-1$ and write
\[
 Q_q[\eta_0,\ldots,\eta_{q-1}]
   =x^q+\sum_{j=0}^{q-1}q_jx^j
\]
for the instance of \Cref{alg:constr-known-2n-1}, with its known powers
regarded as auxiliary coefficients.  Then
\[
 q_j=\eta_j+f_j(\eta_{j+1},\ldots,\eta_{q-1})
 \qquad(0\le j<q).
 \tag{Q-tri}
\]
where the coefficients of $f_j$ may depend polynomially on the supplied
known-power coefficients (and on no active parameter $\eta_i$ with $i\le j$).
In particular its coefficient map is unitriangular from high degree to low
degree.
\end{lemma}
\begin{proof}
We spell out the slot invariant, since mere decodability in
\Cref{lem:fill-correctness} would not by itself imply (Q-tri).

As in the preceding mutual induction, let $\mathsf S_l$ denote the slot
assertion for a level-$l$ fill call, and let $\mathsf U_s$ denote (Q-tri) for
$Q_{2^s-1}$.  The direct formulas give $\mathsf U_1,\mathsf U_2$, and the
argument below proves
\[
  (\mathsf S_1,\ldots,\mathsf S_{l-1};
    \mathsf U_1,\ldots,\mathsf U_l)\Longrightarrow\mathsf S_l,
  \qquad
  (\mathsf S_l;\mathsf U_1,\ldots,\mathsf U_l)
       \Longrightarrow\mathsf U_{l+2}.
  \tag{slot-order}
\]
Thus the order
$\mathsf S_1,\mathsf U_3,\mathsf S_2,\mathsf U_4,\ldots$ is well founded.

Consider a level-$l$ fill call, put $D=2^l$, and let its input pair have
degree $n$.  Suppose a parameter $z$ of the input pair has a descending
affine pivot of slope $\lambda$ in row $g$ of its combined polynomial, and
that its two component coefficients obey the compatibility cutoffs at that
row.  A simultaneous induction on the fill level gives the following three
coefficient bands in the output $A_D$:
\begin{enumerate}
  \item the fill parameters $\alpha_0,\ldots,\alpha_{2D-3}$ have unit
        pivots in rows $0,\ldots,2D-3$, respectively;
  \item the input pivot for $z$ occurs in row $2D-2+g$, still with slope
        $\lambda$;
  \item for $0\le t\le D$, the scalar $\beta_t$ has a unit pivot in row
        $n+2D-2-t$.
\end{enumerate}
In every row, the remaining terms involve only quantities having pivots in
strictly higher rows.

Here is the induction proving this slot assertion.  For $D=2$, expand
\[
 (x+\beta_0)\bigl((H_2+\beta_1)S^{(1)}+\alpha_1\bigr)
   +(H_2+\beta_2)S^{(2)}+\alpha_0.
\]
The parameters $\alpha_0,\alpha_1$ occur first in rows $0,1$; multiplication
by the monic quadratic shifts every input pivot by two rows; and
$\beta_2,\beta_1,\beta_0$ occur first in rows $n,n+1,n+2$.  All these
coefficients are~$1$.  A lower-row occurrence cannot affect a higher pivot,
and the compatibility cutoff says that the extra term
$\beta_0A^{(1)}_2$ cannot reintroduce the parameter whose pivot is currently
being read.

For the induction step, first form the intermediate pair in
\Cref{alg:constr-fill}.  Its three disjoint raw bands are
\[
\begin{array}{c|c}
  0,\ldots,D-1&
    \alpha_{D-2},\alpha_{D-1},\ldots,\alpha_{2D-3},\\
  D+g&\text{an input pivot originally in row }g,\\
  n,\ldots,n+D/2-1&
    \beta_D,\beta_{D-1},\ldots,\beta_{D/2+1}.
\end{array}
\tag{fill-raw-slots}
\]
The first band comes from
$xQ_{D-1}+\alpha_{D-2}$, the second from multiplication by monic
degree-$D$ factors, and the third from
$xQ_{D/2-1}+\beta_D$.  The two $Q$-blocks have unit pivots by the
already-proved $\mathsf U_{l-1}$ and $\mathsf U_l$ cases in
(slot-order).  (At
$D=4$ the displayed special branch is checked directly from $Q_3$; it has
the identical parameter slots.)  The leading term in each Cauchy product
preserves the pivot slope, while all other summands contain a higher
coefficient of the relevant factor.

Now apply the level-$(l-1)$ slot assertion to this intermediate pair.  It
shifts every row in (fill-raw-slots) by $D-2$, creates the new bottom band
$\alpha_0,\ldots,\alpha_{D-3}$, and creates the remaining top scalars
$\beta_{D/2},\ldots,\beta_0$.  The resulting rows are exactly the three
bands stated above.  This proves the slot assertion at level~$l$.

We now prove the second implication in (slot-order) by applying the slot
assertion to the known-powers construction.  The claims for $Q_1$ and
$Q_3$ are immediate from their displayed formulas.  For $Q_7$, the raw
combined polynomial of the pair $(H_4+\alpha_3,H_4+\alpha_2)$, after its
known common part is subtracted, is
$x\alpha_3+\alpha_2$.  The $D=2$ slot calculation therefore places
$\alpha_0,\ldots,\alpha_6$ in rows $0,\ldots,6$ with unit slopes.

For $s\ge4$, put $D=2^{s-2}$, so the input pair to the outer $A_D$ call has degree $2D$.
After subtracting its known common power, its combined polynomial is
\[
  xQ_{D-1}+\alpha_{2D-2}.
\]
Thus its raw row zero exposes $\alpha_{2D-2}$, and raw row $j+1$ exposes
$\alpha_{2D-1+j}$ with unit slope for $0\le j\le D-2$, by induction on
the smaller $Q_{D-1}$.  The three output bands are consequently
\[
\begin{array}{c|c}
  0,\ldots,2D-3&\alpha_0,\ldots,\alpha_{2D-3},\\
  2D-2,\ldots,3D-3&\alpha_{2D-2},\ldots,\alpha_{3D-3},\\
  3D-2,\ldots,4D-2&\alpha_{3D-2},\ldots,\alpha_{4D-2}.
\end{array}
\]
In the last line the call passes
$\beta_D,\ldots,\beta_0$ in precisely the displayed increasing
$\alpha$-order, so the row formula $2D+2D-2-t$ gives the indicated slots.
Since $4D-1=2^s-1=q$, every row $j<q$ therefore has the form
$q_j=\eta_j+f_j(\eta_{j+1},\ldots,\eta_{q-1})$, as required.
\end{proof}

\begin{algorithm}[H]
  \caption{Decoder summary for the fill construction $A_{2^l}$}\label{alg:decode-fill}
  \begin{algorithmic}
    \Require level $l\ge 1$, known powers $(H_2,\ldots,H_{2^L})$ with $L\ge l$, and $P=A_{2^l}(\cdot)$ as in \Cref{lem:fill-correctness} (including the cutoff-respecting decoder witnessing compatibility of the input pair on $G\subseteq\rng{n-2^l}$)
    \Ensure the parameter blocks $(\alpha_0,\ldots,\alpha_{2^{l+1}-3},\beta_{2^l},\ldots,\beta_0)$ and the input pair $(S^{(1)}_{2^l},S^{(2)}_{2^l})$

    \If{$l=1$}
      \State Use the explicit compatibility window $G'$ from the base case in \Cref{lem:fill-correctness} and apply \Cref{lem:x-alpha-extraction} to extract $\beta_0$ and derive $(A^{(1)}_2,A^{(2)}_2)$.
      \State Extract $\beta_1,\beta_2$ from the coefficient of $x^n$ in $A^{(1)}_2,A^{(2)}_2$ (using that the relevant top coefficients of $(S^{(1)}_2,S^{(2)}_2)$ are derivable from the auxiliary data by compatibility).
      \State Apply \Cref{lem:monic-division} to recover $S^{(1)}_2,S^{(2)}_2$ and the remainders $\alpha_1,\alpha_0$.
    \ElsIf{$l=2$}
      \State First run the $l=1$ decoder on $P$ to derive $(S^{(1)}_2,S^{(2)}_2)$ and extract $(\alpha_0,\alpha_1,\beta_0,\beta_1,\beta_2)$.
      \State Undo the $l=2$ fill step as in \Cref{lem:fill-correctness}: extract $\beta_3,\beta_4$ from the coefficient of $x^n$ in the equations
      \(S^{(1)}_2=(H_4+\beta_3)S^{(1)}_4+Q_3\) and \(S^{(2)}_2=(H_4+\beta_4)S^{(2)}_4+\alpha_2\),
      then divide by the monic divisors to recover $(S^{(1)}_4,S^{(2)}_4)$ and $(Q_3,\alpha_2)$.
      \State Decode $Q_3$ by the explicit $Q_3$ lemma to extract $(\alpha_3,\alpha_4,\alpha_5)$.
    \Else \Comment{$l\ge 3$}
      \State First run the decoder recursively at level $l-1$ on $P$ to derive the intermediate pair $(S^{(1)}_{2^{l-1}},S^{(2)}_{2^{l-1}})$ and extract the lower blocks $(\alpha_0,\ldots,\alpha_{2^l-3},\beta_{2^{l-1}},\ldots,\beta_0)$.
      \State Form the combined polynomial $\Phi=xS^{(1)}_{2^{l-1}}+S^{(2)}_{2^{l-1}}$ and subtract the known $x^{2^l}$ contribution (coming from monicity of the degree-$(2^l-1)$ $Q_{\mathrm{high}}$ term) to obtain $\Phi_{\mathrm{prod}}$ as in \Cref{lem:fill-correctness}.
      \State Using compatibility on the disjoint window $(2^l+G)\cup(n+\idx{2^{l-1}-1})$, recover the product pair
      \(
        \bigl((H_{2^l}+Q_{\mathrm{low}})S^{(1)}_{2^l},\,(H_{2^l}+\beta_{2^l})S^{(2)}_{2^l}\bigr)
      \)
      and then (by the multiplicativity reconstruction) derive the factors $(H_{2^l}+Q_{\mathrm{low}},H_{2^l}+\beta_{2^l})$ and $(S^{(1)}_{2^l},S^{(2)}_{2^l})$.
      \State Extract $\beta_{2^l}$ from the constant term of $H_{2^l}+\beta_{2^l}$.
      \State Apply \Cref{lem:monic-division} to recover the remainders $Q_{\mathrm{high}}$ and $\alpha_{2^l-2}$ from the defining equations for $S^{(1)}_{2^{l-1}}$ and $S^{(2)}_{2^{l-1}}$.
      \State Decode $Q_{\mathrm{low}}$ and $Q_{\mathrm{high}}$ by the known-powers construction (base cases $Q_3,Q_7$ and otherwise smaller fill levels), extracting the remaining $\beta$- and $\alpha$-blocks.
    \EndIf
  \end{algorithmic}
\end{algorithm}

\begin{algorithm}[H]
  \caption{Decoder summary for $Q_{2^k-1}$ from known powers}\label{alg:decode-Q-2kminus1}
  \begin{algorithmic}
    \Require integer $k\ge 1$, known powers $(H_2,\ldots,H_{2^{k-1}})$, and $Q=Q_{2^k-1}[\alpha_0,\ldots,\alpha_{2^k-2}]$ as produced by \Cref{alg:constr-known-2n-1}; the peeled variant $Q^{\mathrm{peel}}_{2^k-1}$ of \Cref{sec:peeled-Q} is decoded by \Cref{lem:peeled-Q-decodable} instead
    \Ensure the parameter block $(\alpha_0,\ldots,\alpha_{2^k-2})$

    \If{$k=1$}
      \State Output $\alpha_0=[x^0]Q$.
    \ElsIf{$k=2$}
      \State Apply the explicit $Q_3$ decoding lemma to extract $(\alpha_0,\alpha_1,\alpha_2)$ from $(Q,H_2)$.
    \ElsIf{$k=3$}
      \State In \Cref{alg:constr-known-2n-1} we have $Q_7=A_2(S^{(1)}_2,S^{(2)}_2,(x,H_2))$ with $S^{(1)}_2=H_4+\alpha_3$ and $S^{(2)}_2=H_4+\alpha_2$.
      \State Apply the $A_2$ instance of \Cref{alg:decode-fill} (i.e. $l=1$) to recover $(\alpha_0,\alpha_1,\alpha_4,\alpha_5,\alpha_6)$ and derive $(S^{(1)}_2,S^{(2)}_2)$.
      \State Extract $\alpha_2,\alpha_3$ from the constant-term shifts $S^{(2)}_2-H_4$ and $S^{(1)}_2-H_4$.
    \Else \Comment{$k\ge 4$}
      \State In \Cref{alg:constr-known-2n-1} we have
      \(Q_{2^k-1}=A_{2^{k-2}}(S^{(1)}_{2^{k-2}},S^{(2)}_{2^{k-2}},(x,H_2,\ldots,H_{2^{k-2}}))\)
      with
      \(S^{(1)}_{2^{k-2}}=H_{2^{k-1}}+Q_{2^{k-2}-1}\) and \(S^{(2)}_{2^{k-2}}=H_{2^{k-1}}+\alpha_{2^{k-1}-2}\).
      \State Apply \Cref{alg:decode-fill} at level $l=k-2$ to $Q_{2^k-1}$ to extract the parameters used explicitly in that $A_{2^{k-2}}$ call and to derive the input pair $(S^{(1)}_{2^{k-2}},S^{(2)}_{2^{k-2}})$.
      \State Extract the scalar $\alpha_{2^{k-1}-2}$ from $S^{(2)}_{2^{k-2}}-H_{2^{k-1}}$.
      \State Form the lower instance $Q_{2^{k-2}-1}=S^{(1)}_{2^{k-2}}-H_{2^{k-1}}$ and recurse to decode its parameter block.
      \State Combine the extracted blocks to obtain all $\alpha_0,\ldots,\alpha_{2^k-2}$.
    \EndIf
  \end{algorithmic}
\end{algorithm}


\begin{lemma}[Costs of the fill and Mersenne gadgets]\label{lem:fill-Q-count}
    For all $l\ge 2$ and $k\ge 2$:
    \begin{align}
        \mathcal{A}(A^{(1)}_{2^l}) + \mathcal{A}(A^{(2)}_{2^l})
            &= 5(2^{l - 1} + 2^{l - 2}) - 4 \\
        \mathcal{M}(A^{(1)}_{2^l}) + \mathcal{M}(A^{(2)}_{2^l})
            &= 2^l + 2^{l - 1} - 1 \\
        \mathcal{A}(A_{2^l}) &= 5(2^{l - 1} + 2^{l - 2}) - 2 \\
        \mathcal{M}(A_{2^l}) &= 2^l + 2^{l - 1} \\
        \mathcal{A}(Q_{2^k - 1}) &= 5 \cdot 2^{k - 2} - 2 \\
        \mathcal{M}(Q_{2^k - 1}) &= 2^{k - 1} - 1
    \end{align}
\end{lemma}
\begin{proof}
    We count top-level polynomial additions and multiplications (macro-ops),
    with every shared subexpression charged only once.  To make the mutual
    dependence of the two recursions explicit, write
    \[
      \begin{aligned}
      q_k^A&\coloneqq\mathcal A(Q_{2^k-1}),&
      q_k^M&\coloneqq\mathcal M(Q_{2^k-1}),\\
      a_l&\coloneqq\mathcal A(A^{(1)}_{2^l})+\mathcal A(A^{(2)}_{2^l}),&
      m_l&\coloneqq\mathcal M(A^{(1)}_{2^l})+\mathcal M(A^{(2)}_{2^l}).
      \end{aligned}
    \]
    The two algorithms give, for $k\ge4$ and $l\ge3$ respectively,
    \begin{align}
      q_k^A&=q_{k-2}^A+2+\bigl(a_{k-2}+2\bigr),
      &q_k^M&=q_{k-2}^M+\bigl(m_{k-2}+1\bigr),
      \label{eq:count-q-rec}\\
      a_l&=a_{l-1}+4+q_{l-1}^A+q_l^A,
      &m_l&=m_{l-1}+2+q_{l-1}^M+q_l^M.
      \label{eq:count-a-rec}
    \end{align}
    Here the two additions in the first line of \eqref{eq:count-q-rec}
    form the two inputs to the outer fill call; the final product in
    $A_{2^l}$ is accounted for by the $+2$ and $+1$ in that line.

    The direct bases are
    \[
      (q_2^A,q_2^M)=(3,1),\qquad
      (q_3^A,q_3^M)=(8,3),\qquad
      (a_2,m_2)=(11,5).
    \]
    (For $Q_3$ these are the three additions and one product in
    $(x+\alpha_2)(H_2+\alpha_1)+\alpha_0$.  For $Q_7$, the two additions
    forming $H_4+\alpha_3,H_4+\alpha_2$ together with the six additions and
    three products in $A_2$ give $(8,3)$.)

    We now prove all four closed forms simultaneously.  At stage $r=3$,
    \eqref{eq:count-a-rec} and the bases give
    $(a_3,m_3)=(26,11)$; the $Q_3,Q_7$ bases are already known.  For
    $r\ge4$, assume the formulas for every index $<r$.  First use
    \eqref{eq:count-q-rec}, whose right-hand side involves only index
    $r-2$, to obtain
    \begin{align}
      q_r^A
        &=(5\cdot2^{r-4}-2)+2+
          \bigl(5(2^{r-3}+2^{r-4})-2\bigr)
          =5\cdot2^{r-2}-2,\\
      q_r^M
        &=(2^{r-3}-1)+(2^{r-2}+2^{r-3})
          =2^{r-1}-1.
    \end{align}
    We can then use \eqref{eq:count-a-rec} (now including the just-proved
    $q_r^A,q_r^M$) to obtain
    \begin{align}
      a_r
        &=\bigl(5(2^{r-2}+2^{r-3})-4\bigr)+4
          +(5\cdot2^{r-3}-2)+(5\cdot2^{r-2}-2)\\
        &=5(2^{r-1}+2^{r-2})-4,\\
      m_r
        &=\bigl(2^{r-1}+2^{r-2}-1\bigr)+2
          +(2^{r-2}-1)+(2^{r-1}-1)\\
        &=2^r+2^{r-1}-1.
    \end{align}
    Together with the explicit $r=3$ calculation above, this completes the
    simultaneous induction.  Finally, the outer combination
    $A_{2^l}=(x+\beta_0)A^{(1)}_{2^l}+A^{(2)}_{2^l}$ adds two additions and
    one multiplication:
    \[
      \mathcal A(A_{2^l})=a_l+2
        =5(2^{l-1}+2^{l-2})-2,\qquad
      \mathcal M(A_{2^l})=m_l+1
        =2^l+2^{l-1}.
    \]
    The $Q$ formulas above are $q_k^A=5\cdot2^{k-2}-2$ and
    $q_k^M=2^{k-1}-1$, as claimed.
\end{proof}



\subsection{A depth-balanced known-powers gadget}
\label{sec:peeled-Q}

The chains of \Cref{alg:constr-fill,alg:constr-known-2n-1} multiply
sequentially, so the multiplicative \emph{height} of $Q_{2^k-1}$ --- the
maximum number of nonscalar multiplications on an input-to-output path,
studied further in \Cref{sec:addition-count} --- grows quadratically in
$k$ (\Cref{thm:construction-height}).  The following \emph{peeled}
variant computes a family with the same span, the same multiplication and
addition counts, and height exactly $k$ over any tower with
$h(H_{2^i})\le i$ (\Cref{lem:peeled-Q-count});
\Cref{thm:construction-height-peeled} supplies such towers throughout the
construction.  It keeps the bases $Q_1$ and
$Q_3$ of \Cref{alg:constr-known-2n-1} and, for $k\ge3$, sets
\begin{equation}
  Q^{\mathrm{peel}}_{2^k-1}[\gamma,\vec\alpha,\vec\beta]
  \;=\;
  \bigl(H_{2^{k-1}}+\gamma\bigr)\,
    Q^{\mathrm{peel}}_{2^{k-1}-1}[\vec\alpha]
  \;+\;
  Q^{\mathrm{peel}}_{2^{k-1}-1}[\vec\beta],
  \label{eq:peeled-Q}
\end{equation}
where $\vec\alpha$ and $\vec\beta$ are two fresh parameter blocks of length
$2^{k-1}-1$ each.  Both children recurse in parallel; the glue factor
carries the fresh key $\gamma$.  \Cref{fig:peeled-Q} contrasts the two
schedules and pictures the decoder.

\begin{figure}[H]
  \centering
  \resizebox{\textwidth}{!}{% The peeled known-powers gadget: schedule comparison and decoder strip.
\begin{tikzpicture}[x=1cm,y=1cm,>=Latex]
  % Panel (a): sequential chain vs balanced peel.
  \node[fp panel title] at (0,5.0) {(a) Two schedules, one ledger};
  \node[fp known, minimum width=15mm] (q3) at (0.9,4.3) {$Q_3$};
  \node[fp known, minimum width=15mm] (q7) at (2.7,4.3) {$Q_7$};
  \node[fp known, minimum width=17mm] (q15) at (4.6,4.3) {$Q_{15}$};
  \node[fp label] (qd) at (6.2,4.3) {$\cdots$};
  \node[fp known, minimum width=19mm] (qk) at (7.6,4.3) {$Q_{2^k-1}$};
  \draw[fp flow] (q3) -- (q7);
  \draw[fp flow] (q7) -- (q15);
  \draw[fp flow] (q15) -- (qd);
  \draw[fp flow] (qd) -- (qk);
  \node[fp label, anchor=west, text=fpGray] at (8.8,4.3)
    {sequential fill:\\each level waits for the last\\height $\Theta(k^2)$ in the tower};
  \node[fp active, minimum width=42mm] (root) at (2.6,3.3)
    {$\bigl(H_{2^{k-1}}+\gamma\bigr)\cdot W+B$};
  \node[fp recursive, minimum width=26mm] (wch) at (0.9,2.35)
    {$W=Q^{\mathrm{peel}}_{2^{k-1}-1}[\vec\alpha]$};
  \node[fp recursive, minimum width=26mm] (bch) at (4.3,2.35)
    {$B=Q^{\mathrm{peel}}_{2^{k-1}-1}[\vec\beta]$};
  \draw[fp flow] (wch) -- (root);
  \draw[fp flow] (bch) -- (root);
  \node[fp label, anchor=west, text=fpGray] at (6.6,2.8)
    {peeled \eqref{eq:peeled-Q}: both children in parallel\\
     under one keyed glue factor;\\
     height exactly $k$, identical counts};
  % Panel (b): decoder strip.
  \node[fp panel title] at (0,1.75) {(b) Decoder of \Cref{lem:peeled-Q-decodable}};
  \draw[fp decode] (0,1.12) -- node[fp label, above]
    {descending coefficient degree} (13.4,1.12);
  \node[fp active, text width=33mm, anchor=west] (topw) at (0,0.45)
    {top window $=H_{2^{k-1}}\cdot W$;\\divide by the monic $H$};
  \node[fp pivot, text width=26mm, anchor=west] (grow) at (3.65,0.45)
    {$[x^{m}]$ residual $=\gamma+1$;\\one unit pivot};
  \node[fp second, text width=24mm, anchor=west] (brow) at (6.6,0.45)
    {$B=R-\gamma W$\\by subtraction};
  \node[fp recursive, text width=32mm, anchor=west] (rec) at (9.35,0.45)
    {recurse on $W$ and $B$\\(both children monic)};
  \draw[fp dependence] (topw.south east) to[bend right=10] (grow.south west);
  \draw[fp dependence] (grow.south east) to[bend right=10] (brow.south west);
  \draw[fp dependence] (brow.south east) to[bend right=10] (rec.south west);
\end{tikzpicture}
}
  \caption{The peeled known-powers gadget.  (a) The sequential fill chain
  builds $Q_{2^k-1}$ level by level, so each product waits for the previous
  one; the peeled recursion \eqref{eq:peeled-Q} runs its two half-size
  children in parallel under one keyed glue factor, reaching height exactly
  $k$ at the identical multiplication and addition counts.  (b) The decoder:
  the one-degree offset keeps the top window equal to $H_{2^{k-1}}\cdot W$,
  so $W$ peels off by monic division; the residual reveals $\gamma$ at one
  unit pivot because both children are monic; $B$ follows by subtraction,
  and the decoder recurses on both children.}
  \label{fig:peeled-Q}
\end{figure}

\begin{lemma}[The peeled gadget is decodable]
\label{lem:peeled-Q-decodable}
$Q^{\mathrm{peel}}_{2^k-1}$ is decodable given $H_2,\ldots,H_{2^{k-1}}$.
\end{lemma}
\begin{proof}
Write $W$ and $B$ for the two children in \eqref{eq:peeled-Q} and
$m=2^{k-1}-1$, so $\deg W=\deg B=m$ and both are monic by induction.  Since
$\deg(\gamma W+B)=m<2^{k-1}=\deg H_{2^{k-1}}$, division of
$Q^{\mathrm{peel}}_{2^k-1}$ by the known monic $H_{2^{k-1}}$
(\Cref{lem:monic-division}) returns the quotient $W$ and the remainder
$R=\gamma W+B$ exactly.  Both children are monic of degree $m$, so
$[x^m]R=\gamma+1$, which extracts $\gamma$; then $B=R-\gamma W$.  The two
children are decoded by the induction hypothesis, and the bases are the
lemmas above.  Every step is polynomial in the coefficients, and
re-expansion returns the input, so the parameter map is bijective onto the
monic polynomials of degree $2^k-1$.
\end{proof}

\begin{lemma}[Identical ledger]
\label{lem:peeled-Q-count}
$Q^{\mathrm{peel}}_{2^k-1}$ uses exactly $2^{k-1}-1$ multiplications and
$5\,2^{k-2}-2$ additions, the counts of \Cref{lem:fill-Q-count}, and its
height over towers with $h(H_{2^i})=i$ is exactly~$k$ (and at most $k$
whenever $h(H_{2^i})\le i$).
\end{lemma}
\begin{proof}
The recursion \eqref{eq:peeled-Q} costs one product, one addition inside
the glue factor, and one addition for the let-bound value $B$, on top of
two children: $\mathcal M(k)=1+2\mathcal M(k-1)$ and
$\mathcal A(k)=2+2\mathcal A(k-1)$ with the $k=2$ bases $1$ and $3$.  The
paper ledger $q^M_k=2^{k-1}-1$, $q^A_k=5\,2^{k-2}-2$ satisfies the same
recurrences and bases.  For the height,
$h(Q^{\mathrm{peel}}_{2^k-1})
 =1+\max\bigl(h(H_{2^{k-1}}),h(Q^{\mathrm{peel}}_{2^{k-1}-1})\bigr)
 =1+\max(k-1,k-1)=k$ by induction from $h(Q_3)=2$; the monotone bound
follows since lowering an operand height never raises a gate height.
\end{proof}

\begin{theorem}[Peeled schedules have logarithmic height]
\label{thm:construction-height-peeled}
Substitute $Q^{\mathrm{peel}}$ for every instance of
\Cref{alg:constr-known-2n-1} in the complete degree-$n$ construction.
The resulting schedules compute a decodable monic
degree-$n$ family with exactly the multiplication count of
\Cref{thm:construction-count}, an addition count no larger than that of
\Cref{thm:construction-addition-count}, and height at most
\[
  2\lceil\log_2n\rceil+4 \quad (n\text{ odd}),\qquad
  2\lceil\log_2n\rceil+5 \quad (n\text{ even}).
\]
\end{theorem}
\begin{proof}
Decodability: every consumer of a known-powers gadget --- including the
three special cases $15$, $27$, $31$, whose decoders recover their
$Q_7$-block \emph{values} before descending into them --- first recovers the
gadget's value (its coefficient vector) and then invokes the gadget
decoder: each stage-table row first recovers one coefficient of the
gadget's value (using only monicity, via \eqref{eq:monic-Cauchy}), after
which the consecutive parameter block is extracted by the gadget decoder
and discharged by \Cref{lem:triangular-block-concatenation}.  Under this
value-first reading, \Cref{lem:peeled-Q-decodable} substitutes for
\Cref{lem:Q-unitriangular} at every site.  The counts are
\Cref{lem:peeled-Q-count}.  It remains to bound the height by a ledger
induction over the same components as \Cref{thm:construction-height}; the
point of the substitution is that every recurrence below has bounded
increments, whereas the sequential fill's $h(Q_{2^t-1})\approx h(H_{2^{t-1}})+t$ (cf.\ \eqref{eq:fill-height}, \eqref{eq:tower-height}) made the
tower recurrence quadratic.

\emph{Towers and gadgets.}  Write $h(U)$ for the height of a wire.  We
claim, by simultaneous induction on the level,
\[
  h(H_{2^i})=i,
  \qquad
  h\bigl(Q^{\mathrm{peel}}_{2^t-1}\bigr)=t .
\]
Every displayed power step forms $H_{2^{\ell+1}}$ with one product whose
operands are $H_{2^\ell}$ (height $\ell$) and blocks
$H_{2^{\ell-1}}+Q^{\mathrm{peel}}_{2^{\ell-1}-1}$,
$H_{2^{\ell-2}}+Q^{\mathrm{peel}}_{2^{\ell-2}-1}$,
$Q^{\mathrm{peel}}_{2^{\ell-2}-1}$ of height at most $\ell-1$, so
$h(H_{2^{\ell+1}})=\ell+1$; and \eqref{eq:peeled-Q} gives
$h(Q^{\mathrm{peel}}_{2^t-1})
 =1+\max\bigl(h(H_{2^{t-1}}),h(Q^{\mathrm{peel}}_{2^{t-1}-1})\bigr)=t$
from the base $h(Q_3)=2$.

\emph{The $T$ recursion.}  For $T_{k,2^\ell}$ put
$\lambda=\ell+\lceil\log_2k\rceil$ and let $s(k)$ be the number of odd
values exceeding $1$ in the halving chain of $k$, so
$s(k)\le\lceil\log_2k\rceil=\lambda-\ell$.  Then
\[
  h\bigl(T^{(1)}_{k,2^\ell}\bigr),\,h\bigl(T^{(2)}_{k,2^\ell}\bigr)
  \;\le\;\lambda+s(k).
\]
The base $k=1$ is the tower value $H_{2^\lambda}$ at height $\lambda$; an
even step recurses at $(k/2,\ell+1)$ with the same $\lambda$ after one
tower step; an odd step multiplies the recursive pair by factors of height
at most $\ell$ (a tower level plus peeled blocks of level $<\ell$) and adds
$Q^{\mathrm{peel}}_{2^\ell-1}$ of height $\ell$, contributing exactly the
$+1$ that separates $s(k)$ from $s\bigl(\tfrac{k-1}2\bigr)$.

\emph{Fills and the odd-degree family.}  The fill
$A_{2^\ell}(S^{(1)},S^{(2)})$ is a chain of $\ell+1$ products whose factors
carry towers and peeled blocks of level at most $\ell-1$, so
\[
  h(A_{2^\ell})\le\max\bigl(h(S^{(1)}),h(S^{(2)}),\ell\bigr)+\ell+1.
\]
The
family $Q_{2^{l+1}k+2^l-1}$ ($k\ge1$) runs $T_{2k,2^l}$ and finishes with
$A_{2^{l-1}}$ (for $l=1$, with a single $(x+\beta_0)$ extraction), hence
has height at most
\[
  \bigl(\lambda+s(2k)\bigr)+l+1
  \;\le\;
  2\lambda+1
  \;\le\;
  2\lceil\log_2(m+1)\rceil+1
\]
for its degree $m$: here $s(2k)\le\lambda-l$ absorbs the fill length, and
$\lambda=\lceil\log_2(2^l\cdot 2k)\rceil$ with $2^l\cdot2k\le m+1$,
so $\lambda\le\lceil\log_2(m+1)\rceil$ by monotonicity.  The same bound covers the
good-polynomial gadgets: for $\deg\equiv1\pmod4$ they \emph{are} this
family at $l=1$; $\deg\equiv3\pmod8$ adds one product; and the
$8k+7$ fallback synthesizes one keyed $H_8$ (height $\le4$), runs
$T_{k,8}$, and finishes with the constant-length fill $A_4$, all within
$2\lceil\log_2(m+1)\rceil+4$.

\emph{Branches.}  Let $B(L)$ bound the height of every costed pair of
degree $m$ with $\lceil\log_2m\rceil\le L$, and take
$B(L)=2L+3$.  The $4k+1$ family is one extraction over $T_{2k,2}$:
$1+(\lambda+s)\le2\lambda\le2L$.  The $8k+3$ and $8k+7$ steps add one
square-difference on top of odd-degree or good-polynomial blocks of degree
$m\le(n-1)/2$---since odd $n\ge3$ is not a power of two,
$\lceil\log_2(m+1)\rceil\le L-1$, so these blocks have height at most
$2(L-1)+4=2L+2$---and of a recursive pair of degree at most $(n-1)/4$
(height $\le B(L-2)=2L-1$); hence the new pair has height at most
$2L+3=B(L)$.  The special cases $15,27,31$ are finite, use the same peeled
blocks, and satisfy the bound by direct inspection of their displayed
circuits.  The final combination
$P_n=x\cdot T^{(1)}+T^{(2)}$ adds one product:
$B(L)+1=2\lceil\log_2n\rceil+4$.  For even $n$ the lift from degree
$n-1$ adds one further product on top of the odd-degree schedule, giving
$2\lceil\log_2(n-1)\rceil+5\le2\lceil\log_2n\rceil+5$; this is the
constant $2\lceil\log_2 n\rceil+4+[n\text{ even}]$ of the machine-checked
statement \texttt{polynomial\_height} (\texttt{FastPoly/HeightFinal.lean}),
which covers every $n\ge3$.
\end{proof}

The same peel extends to every odd degree: with
$h=2^{\lfloor\log_2d\rfloor}$, the three-child node
\begin{equation}
  Q_d=(H_h+U)\,W+B,
  \qquad U=Q_{2h-d}\ \text{inside the factor},\quad W=B=Q_{d-h},
  \label{eq:odd-peel}
\end{equation}
uses exactly $(d-1)/2$ multiplications for every odd $d$ at the
known-powers interface.  The key-carrying glue factor is essential: a
balanced product of two \emph{free} half-degree blocks is a factoring
problem, while symmetric square-difference shells, though decodable, pay
about $\tfrac43$ of the multiplications and $2.2\times$ the additions.
Every pivot of the peeled decoders is a monic leading coefficient, so ---
unlike the square-difference machinery --- the recursion is
characteristic-independent
(\texttt{tools/balanced\_gadgets.py} verifies both round-trip directions
up to degree $255$ over $\mathrm{GF}(2^{61}-1)$, $\mathrm{GF}(1009)$,
$\mathbb Q$ and $\mathrm{GF}(2)$).

\begin{remark}
Empirically, over every odd $3\le n\le4095$ (checked by
\texttt{tools/polychain.py} \texttt{--peeled}, which replays the full
encode/decode round trip and the per-degree count assertions):
the maximal height of the peeled schedules
is exactly $2\lceil\log_2n\rceil-2$, attained at $n=2^L-1$ (the halving
chain with the most odd steps) and growing by exactly $+2$ per doubling;
the share-aware addition count is strictly below the sequential
schedule's from $n\approx250$ on (e.g.\ $5593$ against $5895$ at
$n=4093$); the odd-degree families additionally route through the
three-child peel \eqref{eq:odd-peel} wherever enough powers are threaded.
Machine-checked (Lean): the master statement
(\texttt{odd\_realizable\_pairs}) carries, alongside decodability and
the exact multiplication count, the height conjuncts
$h(T^{(1)}),h(T^{(2)})\le2\lceil\log_2n\rceil+3$ and $h(H_2)\le1$,
$h(H_4)\le2$ for the very circuit it exhibits --- the plain
$Q^{\mathrm{peel}}$ substitution of this theorem --- and
\texttt{FastPoly/HeightFinal.lean} closes the ledger end to end:
\texttt{odd\_polynomial\_height} gives $h(P_n)\le2\lceil\log_2n\rceil+4$
for odd $n\ge3$, and \texttt{polynomial\_height} gives
$2\lceil\log_2n\rceil+4+[n\text{ even}]$ for every $n\ge3$, each for one
fixed program with exactly $\lfloor n/2\rfloor+1$ products, exactly the
ledger bound above.
\end{remark}

